% ========================== 第十二节 附件 ====================================
% 备查文件目录及查阅时间、地点；其后为「附件一、××」系列——2025 年招股书
% 精简改革后的申报实践（对照 2026 年真实申报稿）：承诺全文、权利清单等
% 细节内容自正文移入附件。模板选录 6 个必要附件；未选录者及理由见文末注释。
\gssection{附件}

\subsection{备查文件}

（一）发行保荐书；

（二）财务报表及审计报告；

（三）内部控制鉴证报告；

（四）经注册会计师核验的非经常性损益明细表；

（五）法律意见书；

（六）公司章程及公司章程（草案）；

（七）中国证券监督管理委员会同意注册的文件；

（八）其他与本次发行有关的重要文件。

\subsection{查阅时间及地点}

自本招股说明书刊登之日起，投资者可在工作日上午9:00—11:30、下午13:30—17:00，
于以下地点查阅上述备查文件：

（一）发行人：\gsCompany 董事会办公室（\gsRegAddr）；

（二）保荐人（主承销商）：\gsSponsor（\gsSponsorAddr）。

投资者亦可在上海证券交易所网站（www.sse.com.cn）查阅本招股说明书全文及上述
备查文件。

% ============================================================================
\gsannex{附件一}{落实投资者关系管理相关规定的安排、股东投票机制建立情况}

\subsection{投资者关系的主要安排}

公司制定了《投资者关系管理制度》，指定董事会办公室为投资者关系管理部门、董
事会秘书为负责人，通过业绩说明会、上交所投资者互动平台、投资者热线与电子信
箱（ir@example.com）等渠道与投资者沟通，保障投资者的知情权、参与权与监督权。
公司将在上市后持续做好投资者关系管理工作，充分听取中小投资者的意见与诉求。

\subsection{股东投票机制的建立情况}

《公司章程（草案）》及《股东会议事规则》规定：股东会采取现场投票与网络投票
相结合的方式召开，为股东参会行权提供便利；审议影响中小投资者利益的重大事项
时，对中小投资者表决单独计票并披露；选举两名以上董事时实行累积投票制；股东
会审议关联交易事项时，关联股东回避表决。

% ============================================================================
\gsannex{附件二}{与投资者保护相关的承诺}

\subsection{关于股份限售安排、自愿锁定股份的承诺}

控股股东示例控股、实际控制人林××承诺：自公司股票上市之日起36个月内，不转
让或者委托他人管理其直接和间接持有的公司本次发行前已发行的股份，也不由公司
回购该部分股份；公司上市后6个月内如公司股票连续20个交易日的收盘价均低于发
行价，或者上市后6个月期末收盘价低于发行价，所持股份的锁定期限自动延长6个月。

担任公司董事、高级管理人员的股东承诺：除前述锁定期外，在任职期间每年转让的
股份不超过其所持公司股份总数的25\%，离职后半年内不转让其所持公司股份。公司
其他股东按照《公司法》及上交所相关规则履行股份锁定义务。

\subsection{股东持股及减持意向的承诺}

控股股东、实际控制人承诺：锁定期届满后减持公司股份的，将遵守中国证监会及上
交所关于股份减持的相关规定，提前依规披露减持计划；锁定期届满后两年内减持的，
减持价格不低于本次发行的发行价（因除权除息等事项调整后）。

\subsection{稳定股价的措施和承诺}

公司上市后三年内，若公司股票连续20个交易日的收盘价均低于公司最近一期经审计
的每股净资产（因除权除息等事项调整后），公司及相关主体将按照``公司回购股份
——控股股东、实际控制人增持——董事（不含独立董事）、高级管理人员增持''的
顺序启动稳定股价措施，具体以股东会审议通过的稳定股价预案为准。

\subsection{对欺诈发行上市的股份购回承诺}

公司及控股股东、实际控制人承诺：若公司不符合发行上市条件，以欺骗手段骗取发
行注册并已经发行上市的，将依法购回本次公开发行的全部新股。

\subsection{填补被摊薄即期回报的措施及承诺}

公司将通过加快募投项目实施进度、提升经营管理效率、完善利润分配政策等措施填
补本次发行摊薄的即期回报。公司董事、高级管理人员承诺：忠实、勤勉地履行职责，
维护公司和全体股东的合法权益；不无偿或以不公平条件向其他单位或者个人输送利
益，不损害公司利益；将个人薪酬制度及拟公布的股权激励行权条件与公司填补回报
措施的执行情况相挂钩。

\subsection{利润分配政策的承诺}

公司承诺：上市后严格执行《公司章程（草案）》确定的利润分配政策，保持利润分
配政策的连续性和稳定性，每年以现金方式分配的利润不少于当年实现的可供分配利
润的10\%，最近三年以现金方式累计分配的利润不少于最近三年实现的年均可分配利
润的30\%。

\subsection{依法承担赔偿责任的承诺}

公司及其控股股东、实际控制人、董事、高级管理人员承诺：因公司招股说明书及其
他信息披露资料有虚假记载、误导性陈述或者重大遗漏，致使投资者在证券发行和交
易中遭受损失的，将依法赔偿投资者损失。保荐人示例证券承诺：因其为发行人本次
公开发行制作、出具的文件有虚假记载、误导性陈述或者重大遗漏，给投资者造成损
失的，将先行赔偿投资者损失。

\subsection{控股股东、实际控制人避免新增同业竞争的承诺}

控股股东示例控股、实际控制人林××承诺：不以任何形式直接或间接从事与公司主
营业务构成竞争的业务，亦不投资、控制与公司主营业务构成竞争的其他企业；如出
现同业竞争情形，将在合理期限内通过资产重组、业务调整或转让等方式予以消除。

\subsection{关于规范关联交易的承诺}

控股股东示例控股、实际控制人林××承诺：尽量减少并规范与公司之间的关联交易；
确有必要发生的关联交易，将遵循平等、自愿、等价、有偿的原则，依法履行关联交
易决策程序及信息披露义务，不利用关联交易损害公司及其他股东的合法权益。

\subsection{在审期间不进行现金分红的承诺}

公司承诺：自本次发行申报受理之日起至完成发行上市（或终止审核）期间，不进行
现金分红；滚存未分配利润由本次发行完成后的新老股东按持股比例共同享有。

\subsection{未履行承诺的约束措施}

公司及相关责任主体承诺：如未能履行本附件所列各项承诺，将及时披露未能履行的
具体原因并向股东和社会公众投资者道歉；给投资者造成损失的，依法赔偿投资者损
失；并接受暂不领取（或降低）薪酬或津贴、暂停转让所持公司股份等约束措施，直
至相关承诺履行完毕或者相应补救措施实施完毕。

% 正式申报文件中通常还包括「关于股东信息披露的专项承诺」「中介机构的承诺」
% 等，按上述格式补写即可。

% ============================================================================
\gsannex{附件三}{发行人及其子公司拥有的专利权清单}

\subsection{已授权的中国境内发明专利（节选）}

\begin{gstab}
\begin{longtable}{|C{10mm}|L{52mm}|L{34mm}|C{20mm}|C{16mm}|}
\hline
\thA{序号} & \thB{专利名称} & \thB{专利号} & \thB{申请日} & \thB{取得方式} \\ \hline
\endfirsthead
\hline
\thA{序号} & \thB{专利名称} & \thB{专利号} & \thB{申请日} & \thB{取得方式} \\ \hline
\endhead
1 & 一种大规模分布式模型训练的梯度同步方法及系统 & ZL202×1××××××.× & 202×年×月 & 原始取得 \\ \hline
2 & 一种混合专家模型的动态路由方法及装置 & ZL202×1××××××.× & 202×年×月 & 原始取得 \\ \hline
3 & 一种多模态数据统一表征的预训练方法 & ZL202×1××××××.× & 202×年×月 & 原始取得 \\ \hline
4 & 一种基于人类反馈的模型偏好对齐训练方法 & ZL202×1××××××.× & 202×年×月 & 原始取得 \\ \hline
5 & 一种大语言模型推理的量化压缩方法及系统 & ZL202×1××××××.× & 202×年×月 & 原始取得 \\ \hline
6 & 一种智能体任务规划与工具调用的编排方法 & ZL202×1××××××.× & 202×年×月 & 原始取得 \\ \hline
7 & 一种生成内容安全过滤与幻觉检测方法 & ZL202×1××××××.× & 202×年×月 & 原始取得 \\ \hline
8 & 一种训练语料去重与质量评估方法及装置 & ZL202×1××××××.× & 202×年×月 & 原始取得 \\ \hline
\end{longtable}
\end{gstab}
\tabnote{以上为示例节选；截至2025年12月31日公司共拥有境内发明专利186项，正
式申报文件中应逐项完整列示，境外专利、实用新型及外观设计专利清单按同一格式
另列。}

% ============================================================================
\gsannex{附件四}{发行人及其子公司拥有的软件著作权清单}

\begin{gstab}
\begin{longtable}{|C{10mm}|L{62mm}|L{36mm}|C{22mm}|}
\hline
\thA{序号} & \thB{软件名称} & \thB{登记号} & \thB{取得方式} \\ \hline
\endfirsthead
\hline
\thA{序号} & \thB{软件名称} & \thB{登记号} & \thB{取得方式} \\ \hline
\endhead
1 & 示例大模型分布式训练平台软件 & 202×SR××××××× & 原始取得 \\ \hline
2 & 示例大模型推理服务引擎软件 & 202×SR××××××× & 原始取得 \\ \hline
3 & 示例开放平台（MaaS）服务系统软件 & 202×SR××××××× & 原始取得 \\ \hline
4 & 示例智能体编排平台软件 & 202×SR××××××× & 原始取得 \\ \hline
5 & 示例训练语料管理与标注系统软件 & 202×SR××××××× & 原始取得 \\ \hline
6 & 示例智能助手应用软件 & 202×SR××××××× & 原始取得 \\ \hline
\end{longtable}
\end{gstab}
\tabnote{以上为示例节选；截至2025年12月31日公司共拥有软件著作权120项，正式
申报文件中应逐项完整列示。}

% ============================================================================
\gsannex{附件五}{发行人及其子公司拥有的主要业务许可及备案清单}

\begin{gstab}
\begin{longtable}{|C{10mm}|L{48mm}|L{30mm}|C{24mm}|C{20mm}|}
\hline
\thA{序号} & \thB{许可／备案名称} & \thB{编号} & \thB{持有主体} & \thB{有效期} \\ \hline
\endfirsthead
\hline
\thA{序号} & \thB{许可／备案名称} & \thB{编号} & \thB{持有主体} & \thB{有效期} \\ \hline
\endhead
1 & 增值电信业务经营许可证 & ×ICP证×××××号 & 发行人 & 202×年×月至202×年×月 \\ \hline
2 & 生成式人工智能服务备案 & ××××××××× & 发行人 & 长期 \\ \hline
3 & 互联网信息服务算法备案 & 网信算备×××××××××××号 & 发行人 & 长期 \\ \hline
4 & 高新技术企业证书 & GR×××××××××× & 发行人 & 202×年×月至202×年×月 \\ \hline
5 & 高新技术企业证书 & GR×××××××××× & 示例智算 & 202×年×月至202×年×月 \\ \hline
\end{longtable}
\end{gstab}

% ============================================================================
\gsannex{附件六}{申报前十二个月新增股东基本信息}

本次申报前十二个月内，公司未发生增资扩股事项，亦无新增股东；公司股东之间未
发生导致公司控制权变化的股份转让。

% ----------------------------------------------------------------------------
% 未选录的附件及理由（对照 2026 年真实申报稿的附件一至十二）：
%   · 三会制度运行情况说明、审计委员会及专门委员会设置情况说明
%     ——本模板第八节「公司治理与独立性」一、二已覆盖；如监管要求单列，
%       将第八节相应内容抽出扩写为附件即可（\gsannex 同格式）。
%   · 集成电路布图设计专有权清单、商标清单、版权清单
%     ——视发行人实际持有情况增删，格式同附件三/附件四。
%   · 募集资金具体运用情况
%     ——本模板第七节保留了完整正文（准则第 57 号要求的章节）；部分 2026 年
%       申报稿将可行性分析等细节移入该附件，如需可将第七节内容拆分扩写。
% ----------------------------------------------------------------------------
